\subsection{مبانی نظری و فرمول‌بندی ریاضی}
فرآیند فشرده‌سازی در ماژول \lr{SAES-SVD v2} در سه مرحله ریاضی پیوسته تدوین می‌گردد:

\begin{enumerate}
    \item \textbf{تجزیه مقادیر منفرد پایه و سنجش چگالی طیفی:}
    ماتریس وزن اولیه $\mathbf{W}_0 \in \mathbb{R}^{d_{\text{out}} \times d_{\text{in}}}$ با رتبه هدف $r$ تحت تقطیع اولیه قرار می‌گیرد:
    \begin{equation}
    \mathbf{W}_0 \approx \mathbf{U}_r \mathbf{\Sigma}_r \mathbf{V}_r^\top = \mathbf{A}_0 \mathbf{B}_0
    \end{equation}
    که در آن فاکتورهای اولیه به صورت $\mathbf{A}_0 = \mathbf{U}_r \mathbf{\Sigma}_r^{1/2} \in \mathbb{R}^{d_{\text{out}} \times r}$ و $\mathbf{B}_0 = \mathbf{\Sigma}_r^{1/2} \mathbf{V}_r^\top \in \mathbb{R}^{r \times d_{\text{in}}}$ مقداردهی می‌شوند. چگالی توزیع انرژی برای هر مقدار منفرد $\sigma_k$ بر پایه رابطه زیر سنجیده می‌شود:
    \begin{equation}
    \rho_k = \frac{\sigma_k^2}{\sum_{j=1}^{\min(d_{\text{out}}, d_{\text{in}})} \sigma_j^2}
    \end{equation}

    \item \textbf{تکرارهای فرونشانی خطا بر مبنای کمترین مربعات منظم‌شده:}
    در هر گام تکرار $t \ge 0$، ماتریس خطای تقریب باقی‌مانده محاسبه می‌گردد:
    \begin{equation}
    \mathbf{E}_t = \mathbf{W}_0 - \mathbf{A}_t \mathbf{B}_t
    \end{equation}
    سپس جهت فرونشاندن $\mathbf{E}_t$ بدون افزایش رتبه فاکتورها، ماتریس‌های $\mathbf{A}$ و $\mathbf{B}$ به ترتیب با حل فرم بسته کمترین مربعات تیخونوف به‌روزرسانی می‌شوند:
    \begin{equation}
    \mathbf{A}_{t+1} = \mathbf{A}_t + \mathbf{E}_t \mathbf{B}_t^\top \left(\mathbf{B}_t \mathbf{B}_t^\top + \epsilon \mathbf{I}_r\right)^{-1}
    \end{equation}
    \begin{equation}
    \mathbf{B}_{t+1} = \mathbf{B}_t + \left(\mathbf{A}_{t+1}^\top \mathbf{A}_{t+1} + \epsilon \mathbf{I}_r\right)^{-1} \mathbf{A}_{t+1}^\top \mathbf{E}_t
    \end{equation}
    که در آن $\epsilon > 0$ ضریب تثبیت‌کننده عددی جهت تضمین خوش‌شرطی وارون ماتریس است.

    \item \textbf{همگرایی و بازسازی خروجی:}
    پس از طی $T$ گام فرونشانی، لایه خطی استاندارد به یک ساختار ترکیبی \lr{SVDLinear} تبدیل می‌گردد که خروجی را به صورت مستقیم محاسبه می‌کند:
    \begin{equation}
    \mathbf{y} = \mathbf{x} \mathbf{B}_T^\top \mathbf{A}_T^\top + \mathbf{b}
    \end{equation}
\end{enumerate}


\paragraph{تحلیل طیفی و نامساوی‌های پایداری خطا}
برای بررسی میزان دقت تقریب ماتریس وزن یا نگاشت توجه در ماژول \lr{saes\_svd\_v2}، تفاضل ماتریس فاکتورگیری‌شده از نگاشت متراکم اصلی بر اساس نامساوی هافدینگ-فون‌نیومان ارزیابی می‌گردد. فرض کنید $\mathbf{W} \in \mathbb{R}^{d_{\text{in}} \times d_{\text{out}}}$ ماتریس وزن هدف باشد که توسط تقریب کم‌مرتبه یا پراکنده با رتبه موثر $r$ جایگزین شده است. کران بالای خطای بازسازی بر مبنای نرم فروبنیوس به صورت زیر تصریح می‌شود:
\begin{equation}
\label{eq:spectral_bound_saes_svd_v2}
\| \mathbf{W} - \mathbf{W}_{r} \|_F = \sqrt{\sum_{i=r+1}^{\min(d_{\text{in}}, d_{\text{out}})} \sigma_i^2(\mathbf{W})} \le \epsilon \| \mathbf{W} \|_F
\end{equation}
که در آن $\sigma_i(\mathbf{W})$ بیانگر $i$-امین مقدار منفرد مرتب‌شده ماتریس به صورت نزولی است. نسبت حفظ انرژی طیفی $\eta_r = \frac{\sum_{i=1}^r \sigma_i^2}{\sum_{i=1}^{d} \sigma_i^2}$ تضمین می‌کند که با انتخاب آستانه $\tau \ge 0.90$، بیش از نود درصد از پراکندگی واریانس داده‌ها حفظ شده و افت سیگنال در انتشار گرادیان رو به عقب به حداقل ممکن کاهش می‌یابد.

\paragraph{تحلیل مجانبی و پیچیدگی اشتراک حافظه سخت‌افزاری}
از منظر پیچیدگی محاسباتی و انتقال داده‌ها در سلسومراتب حافظه (حافظه سراسری \lr{HBM} به حافظه محلی \lr{SRAM})، عملیات ضرب ماتریسی در این مکانیزم نیازمند تخصیص بافر موقت با ابعاد فاکتورگیری‌شده است. برای پردازش دنباله‌ای به طول $N$ با بعد مخفی $D$، بار محاسباتی هسته از فرم درجه دو به رابطه زیر تقلیل می‌یابد:
\begin{equation}
\label{eq:complexity_saes_svd_v2}
\mathcal{C}_{\text{saes\_svd\_v2}} = \mathcal{O}\left( B \cdot N \cdot D \cdot r + B \cdot H \cdot N \cdot d_k \right) \ll \mathcal{O}\left( B \cdot H \cdot N^2 \cdot d_k \right)
\end{equation}
این بهبود مرتبه پیچیدگی، با کاهش چشمگیر حجم تبادل تانسورها در گذرگاه‌های ارتباطی کارت گرافیک، مانع از اشباع پهنای باند حافظه (پدیده \lr{Memory-Bound}) در فرآیند استنتاج طولانی‌مدت می‌گردد.
